\documentclass[11pt,a4paper]{article}
\usepackage{fontspec}
\setmainfont{TeXGyrePagella}
\usepackage[french]{babel}
\usepackage{graphicx,tikz,xcolor,geometry}
\usetikzlibrary{calc}
\geometry{margin=0pt}
\pagestyle{empty}
\definecolor{bleufonce}{HTML}{1B3A6B}
\definecolor{bleugris}{HTML}{5C7190}
\newcommand{\pos}[4]{\node[anchor=north west,inner sep=0pt,outer sep=0pt,text width=#3,align=left] at ($(current page.north west)+(#1,-#2)$) {#4};}
\newcommand{\sect}[1]{{\color{bleufonce}\bfseries\fontsize{17}{19}\selectfont #1}}
\newcommand{\step}[1]{{\color{bleufonce}\bfseries\fontsize{12.5}{14.5}\selectfont #1}}
\newcommand{\smallfont}{\fontsize{10}{12.2}\selectfont}
\newcommand{\shot}[4]{\node[anchor=north west,inner sep=0pt,outer sep=0pt] at ($(current page.north west)+(#1,-#2)$) {\includegraphics[width=#3]{#4}};}
\begin{document}
\begin{tikzpicture}[remember picture,overlay]
\draw[bleufonce,line width=.95pt,rounded corners=8mm]
 ($(current page.north west)+(5mm,-5mm)$) rectangle ($(current page.south east)+(-5mm,5mm)$);

\pos{10mm}{13mm}{189mm}{\sect{2. Créer un nouveau projet dans Proteus}}
\pos{10mm}{25mm}{188mm}{{\fontsize{10.5}{13}\selectfont Avant de dessiner un circuit, il faut créer et enregistrer un projet. Les trois étapes ci-dessous sont illustrées par des captures d'écran de \textbf{Proteus 8}.}}

\pos{10mm}{44mm}{188mm}{\step{Étape 1 --- Ouvrir la création d'un projet}}
\shot{14mm}{57mm}{113mm}{assets/01_accueil.png}
\draw[bleufonce!75,line width=.55pt,rounded corners=1mm]
 ($(current page.north west)+(13.6mm,-56.6mm)$) rectangle ($(current page.north west)+(127.4mm,-77.8mm)$);
\pos{135mm}{58mm}{61mm}{{\smallfont Sur la page d'accueil de Proteus, cliquer sur \textbf{New Project}. L'assistant de création s'ouvre.}}
\pos{14mm}{81mm}{184mm}{{\color{bleugris}\fontsize{8.5}{10}\selectfont Figure 2 --- Commande \textbf{New Project} sur la page d'accueil.}}

\pos{10mm}{92mm}{188mm}{\step{Étape 2 --- Nommer et enregistrer le projet}}
\shot{14mm}{106mm}{126mm}{assets/02_assistant.png}
\draw[bleufonce!75,line width=.55pt,rounded corners=1mm]
 ($(current page.north west)+(13.6mm,-105.6mm)$) rectangle ($(current page.north west)+(140.4mm,-138.3mm)$);
\pos{146mm}{108mm}{50mm}{{\fontsize{9.65}{11.8}\selectfont
 \textbf{Name :} nommer le projet, par exemple \emph{TP01\_LED}.\\[1.9mm]
 \textbf{Path :} choisir son dossier.\\[1.9mm]
 Garder \textbf{New Project}.}}
\pos{14mm}{142mm}{95mm}{{\fontsize{9.6}{11.2}\selectfont Valider avec \textbf{Next}, puis poursuivre l'assistant.}}
\shot{120mm}{141mm}{28mm}{assets/02_bouton_suivant.png}
\pos{14mm}{153mm}{184mm}{{\color{bleugris}\fontsize{8.5}{10}\selectfont Figure 3 --- Champs du projet et bouton de validation.}}

\pos{10mm}{164mm}{188mm}{\step{Étape 3 --- Accéder à la feuille de schéma}}
\shot{14mm}{179mm}{182mm}{assets/03_feuille.png}
\draw[bleufonce!75,line width=.55pt,rounded corners=1mm]
 ($(current page.north west)+(13.6mm,-178.6mm)$) rectangle ($(current page.north west)+(196.4mm,-242.0mm)$);
\pos{14mm}{243mm}{182mm}{{\color{bleugris}\fontsize{8.4}{10}\selectfont Figure 4 --- La feuille de schéma quadrillée de \textbf{Schematic Capture (ISIS)}.}}

\draw[bleufonce,line width=.7pt,rounded corners=2mm]
 ($(current page.north west)+(10mm,-250mm)$) rectangle ($(current page.north west)+(200mm,-260mm)$);
\pos{15mm}{251.3mm}{180mm}{{\fontsize{9.5}{11.4}\selectfont \textcolor{bleufonce}{\bfseries À retenir :} enregistrer le projet, puis utiliser la feuille de schéma pour y placer et relier les composants.}}

\draw[bleufonce!55,line width=.45pt] ($(current page.south west)+(10mm,14.7mm)$)--($(current page.south east)+(-10mm,14.7mm)$);
\node[anchor=south west,inner sep=0pt,text=bleufonce!80] at ($(current page.south west)+(10mm,7.4mm)$) {\fontsize{9.5}{11}\selectfont 2026/2027};
\node[anchor=south east,inner sep=0pt,text=bleufonce!80] at ($(current page.south east)+(-10mm,7.4mm)$) {\fontsize{9.5}{11}\selectfont Page 3};
\end{tikzpicture}
\null
\end{document}
